\chapter{The Linux Underneath}\label{cha:linux}

Every emulated machine sits on a real one. On most of them that fact is
carefully hidden; here it is a feature, because the real machine has
compilers on it and the fantasy one is being used to write programs.

This chapter is about the seam: the one character that crosses it, the
two words that use it to compile your programs, and what is on the other
side.

\section{\texttt{!} --- the Linux prompt, from the machine's prompt}
A line beginning with \verb|!| is not the machine's:

\begin{type}
!ls -l
!git status
!nvim notes.txt
\end{type}

\noindent runs that command on the Linux the machine is standing on, with
its output drawn on the machine's screen. A bare \verb|!| with nothing
after it gives you an interactive shell there; leave it the way you leave
any shell.

It is carried on the Tube (Chapter~\ref{cha:tube}), which is the part of
this machine already built to have another processor's console on the
glass: the host's shell is started on the same pty the co-processors use,
JIM renders it, and the ROM's own key loop feeds it. So the colours stay
the machine's, full-screen programs work --- \verb|nvim| runs in the
window --- and \verb|ls --color| lands in the machine's palette. The
terminal type to tell anything that asks is \textbf{ansi}: JIM has
colour, and \texttt{vt100}'s terminfo does not.

\begin{aside}
\textbf{What this means for the container.} The sandboxed flavour
(\pth{linux/podman.sh}) is a wall on purpose: it is handed the display,
the sound, the game controllers and one folder, which the machine sees
as \pth{/MNT/SHARE}, and nothing else. \verb|!| inside it is real, and
reaches the container --- which is a small Debian with the compilers on
it, and not your computer. That is the whole point of running it that
way.
\end{aside}

\section{\texttt{PAS} and \texttt{CC}: compiling on the machine}
The compilers live on the Linux side, so the machine reaches them the
same way you would --- but without you having to remember four options:

\begin{type}
VI HELLO.PAS
PAS HELLO
HELLO
\end{type}

\verb|PAS name| compiles \texttt{name.PAS} \emph{in the directory you are
standing in} with Mad Pascal, and leaves \texttt{name.prg} beside it.
\verb|CC name| does the same for \texttt{name.C} with cc65, using exactly
the rule the machine's own Makefile uses for its programs --- the C output
is byte-for-byte what the repository ships. Intermediates go to a scratch
directory and never appear on your disk; the compiler's own errors appear
on the screen, and its exit status comes back as the result code, so a
script can test it.

That is the loop the machine owns: \emph{edit the source on the machine,
compile it from the machine's prompt, run it on the machine.} The
compiler itself still runs on the Linux beside it, which is why
``self-hosted'' has an ``almost'' in front of it --- but the part that
matters when you are writing a program is all on this side of the seam.

Both are words in the ROM rather than aliases, deliberately: an alias
would live in a \texttt{STARTUP.BAT}, which is per-machine and yours, and
so would never reach a machine you had just built.

\section{What is on the other side}
On a desktop, the other side is your own computer and has whatever you
have put on it. On the stick it is a Debian built for the purpose, and
the list is short because everything on it is there for a reason:

\begin{center}
\small
\begin{tabular}{@{}l >{\raggedright\arraybackslash}p{0.58\linewidth}@{}}
\toprule
cc65, 64tass, nasm & the machine's own toolchain --- it can rebuild itself \\
FPC + Mad Pascal & what \verb|PAS| runs \\
git, neovim & so that work done here is work \\
\texttt{tek40xx} & a Tektronix 4010 terminal, below \\
telnetd & bound to loopback and nothing else \\
NetworkManager & \texttt{nmtui} on another console, for wireless \\
\bottomrule
\end{tabular}
\end{center}

\textbf{The consoles.} Ctrl+Alt+F1 is the machine; Ctrl+Alt+F2 is a
second one, where the Tektronix lives. Quitting the emulator (Shift+Esc,
or F7 $\rightarrow$ Quit) drops you to a Linux shell on the machine's own
console rather than to nothing, and F7 has a \emph{shut down the
computer} row underneath it because on the stick there is no desktop to
go back to. The power button performs a clean shutdown.

\textbf{The internal drive is not there.} The stick's kernel command line
blacklists the storage drivers absolutely, so the laptop's own disk is
not mounted, not visible, and not installable-to. That is not a policy
setting you could turn off by accident; the machine physically cannot see
it. Everything is loaded into RAM at boot, so the stick could be pulled
out --- except that a small persistence partition on it keeps your home
directory, which is where \texttt{k4510.cfg}, your saved states and your
files are. Leave the stick in.

\textbf{Telnet, both directions.} \verb|TELNET 127.0.0.1 23| from the
machine's prompt logs in to the Linux underneath through the front door
rather than through \verb|!| --- useful when you want a session that
survives what the machine is doing. The daemon is bound to the loopback
address only, so there is no network path to it at all. Outbound, the
network is real: \verb|TELNET| reaches a BBS, and \texttt{http://} and
\texttt{tnfs://} work as Chapter~\ref{cha:shell} describes.

\section{A Tektronix beside the machine}
\texttt{tek40xx} (Ian Schofield's, GPL-3) is a Tektronix 4010/4014
storage-tube terminal that is also a telnet client. It is on the stick and
in the container because of what a storage tube is \emph{for}: a PiDP-11
answers on telnet ports, and a machine of that vintage deserves a screen
of that vintage to draw on.

\begin{type}
tek HOST [PORT]
\end{type}

\noindent on the second console --- full screen on a bare console, in a
window under a desktop. And without any PDP-11 at all:

\begin{type}
tekplay NAME|FILE...
\end{type}

\noindent plays Tektronix plot files at a 9600-baud pace, which is the
right speed to watch one being drawn. With no argument it plays every
plot that ships; \texttt{HOME} clears the screen and \texttt{END} quits.
Nothing in the emulator changed for any of this: it is a second program
on the same computer, which is exactly what a second terminal was in
1975.
